\subsection{HALL SETS IN A FREE MAGMA}

Let \(\mathbf{X}\) be a set, \(\mathrm{M}(\mathbf{X})\) the free magma constructed on \(\mathbf{X}\) and \(\mathrm{M}^n (\mathbf{X})\) , where \(n\in \mathbf{N}^*\) , the set of elements of \(\mathrm{M}(\mathbf{X})\) of length \(n\) (no. 1). If \(w\in \mathrm{M}(\mathbf{X})\) and \(l(w)\geqslant 2\) , let \(\alpha (w)\) and \(\beta (w)\) denote the elements of \(\mathrm{M}(\mathbf{X})\) determined by the relation \(w = \alpha (w)\beta (w)\) ; then \(l(\alpha (w)) < l(w)\) , \(l(\beta (w)) < l(w)\) . Finally, for \(u,v\) in \(\mathrm{M}(\mathbf{X})\) , let \(u^n v\) denote the element defined by induction on the integer \(m\geqslant 0\) by \(u^{0}v = v\) and \(u^{m + 1}v = u(u^{m}v)\) .

DEFINITION 2. A Hall set relative to \(\mathbf{X}\) is any totally ordered subset \(\mathbf{H}\) of \(\mathbf{M}(\mathbf{X})\) satisfying the following conditions:

\begin{enumerate}
\def\labelenumi{(\Alph{enumi})}
\item
  If \(u \in \mathbf{H}\) , \(v \in \mathbf{H}\) and \(l(u) < l(v)\) , then \(u < v\) .
\item
  \(\mathbf{X} \subset \mathbf{H}\) and \(\mathrm{H} \cap \mathrm{M}^2(\mathbf{X})\) consists of the products \(xy\) with \(x, y\) in \(\mathbf{X}\) and \(x < y\) .
\item
  An element \(w\) of \(\mathbf{M}(\mathbf{X})\) of length \(\geqslant 3\) belongs to \(\mathbf{H}\) if and only if it is of the form \(a(bc)\) with \(a, b, c\) in \(\mathbf{H}, bc \in \mathbf{H}, b \leqslant a < bc\) and \(b < c\) .
\end{enumerate}

PROPOSITION 11. There exists a Hall set relative to X.

We shall construct by induction on the integer \(n\) sets \(\mathbf{H}_n \subset \mathbf{M}^n(\mathbf{X})\) and a total ordering on these sets:

\begin{enumerate}
\def\labelenumi{(\alph{enumi})}
\item
  We write \(\mathbf{H}_1 = \mathbf{X}\) and give it a total ordering.
\item
  The set \(\mathbf{H}_2\) consists of the products \(xy\) with \(x, y\) in \(\mathbf{X}\) and \(x < y\) . We give it a total ordering.
\item
  Let \(n \geqslant 3\) be such that the totally ordered sets \(\mathrm{H}_1, \ldots, \mathrm{H}_{n-1}\) are already defined. The set \(\mathrm{H}_{n-1}' = \mathrm{H}_1 \cup \dots \cup \mathrm{H}_{n-1}\) has a total ordering which induces the given relations on \(\mathrm{H}_1, \ldots, \mathrm{H}_{n-1}\) and is such that \(w < w'\) if \(l(w) < l(w')\) . We define \(\mathrm{H}_n\) to be the set of products \(a(bc) \in \mathrm{M}^n(\mathbf{X})\) with \(a, b, c\) in \(\mathrm{H}_{n-1}'\) satisfying the relations \(bc \in \mathrm{H}_{n-1}', b \leqslant a < bc, b < c\) and give \(\mathrm{H}_n\) a total ordering.
\end{enumerate}

We write \(\mathbf{H} = \bigcup_{n\geqslant 1}\mathbf{H}_n\) ; we give H the total ordering defined thus: \(w\leqslant w'\) if and only if \(l(w) < l(w')\) or \(l(w) = l(w') = n\) and \(w\leqslant w'\) in the set \(\mathbf{H}_n\) . It is immediate that H is a Hall set relative to X.

For every subset S of X, we identify the free magma M(S) with its canonical image in M(X).

PROPOSITION 12. Let H be a Hall set relative to X and let x, y be in X.

\begin{enumerate}
\def\labelenumi{(\alph{enumi})}
\item
  \(H \cap M(\{x\}) = \{x\}\) .
\item
  Suppose that \(x < y\) and let \(d_y\) be the homomorphism of \(\mathbf{M}(\mathbf{X})\) into \(\mathbf{N}\) such that \(d_y(y) = 1\) and \(d_y(z) = 0\) for \(z \in \mathbf{X}, z \neq y\) . The set of elements \(w \in \mathbf{H} \cap \mathbf{M}(\{x, y\})\) such that \(d_y(w) = 1\) consists of the elements \(x^n y\) with \(n\) an integer \(\geqslant 0\) .
\end{enumerate}

By Definition 2 (B), \(x \in \mathrm{H}\) and \(\mathrm{H} \cap \mathrm{M}^2(\{x\}) = \varnothing\) . If \(w \in \mathrm{H} \cap \mathrm{M}(\{x\})\) , where \(n = l(w) \geqslant 3\) , the elements \(\alpha(w)\) and \(\beta(w)\) also belong to \(\mathrm{H} \cap \mathrm{M}(\{x\})\) by Definition 2 (C). It immediately follows by induction on \(n\) that \(\mathrm{H} \cap \mathrm{M}^n(\{x\}) = \varnothing\) for \(n \geqslant 2\) , whence (a).

We now prove (b). By Definition 2 (B), \(y \in \mathsf{H}\) and \(xy \in \mathsf{H}\) . We show by induction on \(n\) that \(x^n y \in \mathsf{H}\) for \(n\) an integer \(\geqslant 2\) . Now \(x^n y = x(x^{n-2}y)\) and the induction hypothesis implies that \(x^{n-2}y \in \mathsf{H}\) . Now \(l(x) < l(x^{n-2}y)\) for \(n > 2\) and \(x < y\) , whence \(x < x^{n-2}y\) in any case; condition (C) of Definition 2 shows that \(x^n y \in \mathsf{H}\) . On the other hand, certainly \(d_y(x^n y) = 1\) . Conversely, let \(w \in \mathsf{H} \cap \mathrm{M}(\{x, y\})\) , with \(d_y(w) = 1\) . If \(l(w) = 1\) , then \(w = y\) ; if \(l(w) = 2\) , then \(w = xy\) by Definition 2 (B). If \(l(w) \geqslant 3\) , then \(w = a(bc)\) , with \(a, b, c, bc\) in \(\mathsf{H} \cap \mathrm{M}(\{x, y\})\) (Definition 2 (C)). \(d_y(bc) = 0\) is impossible, since this would imply \(bc \in \mathrm{M}(\{x\})\) , which is impossible by (a). Hence \(d_y(bc) = 1\) and \(d_y(a) = 0\) , whence \(a = x\) by (a). It follows immediately by induction on \(n = l(w)\) that \(w = x^{n-2}y\) , which completes the proof of (b).

COROLLARY. If \(\operatorname{Card} X \geqslant 2\) , then \(H \cap M^n(X) \neq \varnothing\) for every integer \(n \geqslant 1\) .

PROPOSITION 13. Let X be a finite set with at least two elements. Let H denote a Hall set relative to X. Then there exist a strictly increasing bijection \(p \mapsto w_p\) of N onto H and a sequence \((\mathbf{P}_p)_{p \in \mathbf{N}}\) of subsets of H with the following properties:

\begin{enumerate}
\def\labelenumi{(\alph{enumi})}
\item
  \(P_{0} = X.\)
\item
  For every integer \(p \geqslant 0\) , \(w_{p} \in \mathbf{P}_{p}\) .
\item
  For every integer \(n \geqslant 1\) , there exists an integer \(p(n)\) such that every element of \(\mathbf{P}_p\) is of length \(> n\) for all \(p \geqslant p(n)\) .
\item
  For every integer \(p \geqslant 0\) , the set \(\mathrm{P}_{p+1}\) consists of the elements of the form \(w_p^i w\) , where \(i \geqslant 0, w \in \mathrm{P}_p\) and \(w \neq w_p\) .
\end{enumerate}

As X is finite, each of the sets \(\mathrm{M}^{n}(\mathrm{X})\) is finite. Let \(H_{n}=H\cap M^{n}(\mathbf{X})\) for all \(n\geqslant1\) . The Corollary to Proposition 12 shows that the finite set \(H_{n}\) is non-empty. Let \(u_{n}\) be the cardinal of \(H_{n}\) ; let \(v_{0}=0\) and \(v_{n}=u_{1}+\cdots+u_{n}\) for \(n\geqslant1\) . As \(H_{n}\) is a totally ordered finite set, there exists a strictly increasing bijection \(p\mapsto w_{p}\) of the interval \([v_{n-1},v_{n}-1]\) of N onto \(H_{n}\) . It is immediate that \(p\mapsto w_{p}\) is a strictly increasing bijection of N onto H.

Let \(\mathbf{P}_0 = \mathbf{X}\) and for every integer \(p \geqslant 1\) let \(\mathbf{P}_p\) be the set of elements \(w\) of \(\mathbf{H}\) such that \(w \geqslant w_p\) , and, either \(w \in \mathbf{X}\) , or \(\alpha(w) < w_p\) (note that if \(w\) is of length \(\geqslant2\) the relation \(w\in H\) implies \(\alpha(w)\in H\) by condition (C) of Definition 2). Then \(w_{p}\in P_{p}\) ; this is clear if \(w_{p}\in X\) and follows from the inequality \(l(\alpha(w_{p}))<l(w_{p})\) and condition (A) of Definition 2 when \(w_{p}\notin X\) .

Hence conditions (a) and (b) are satisfied.

Let \(n\) be an integer \(\geqslant 1\) and let \(p \geqslant v_n\) . For all \(w \in \mathrm{P}_p\) , \(l(w) \geqslant l(w_p) > n\) by the very definition of the mapping \(p \mapsto w_p\) . This establishes (c).

We now show that every element of the form \(u = w_{p}^{i}w\) with \(i \geqslant 0\) , \(w \in \mathrm{P}_{p}\) and \(w \neq w_{p}\) belongs to \(\mathrm{P}_{p+1}\) . If \(i \neq 0\) , then \(l(u) > l(w_{p})\) , whence \(u > w_{p}\) and \(u \geqslant w_{p+1}\) ; then \(u \notin X\) and \(\alpha(u) = w_{p} < w_{p+1}\) , whence \(u \in \mathrm{P}_{p+1}\) . If \(i = 0\) , then \(u \in \mathrm{P}_{p}\) and \(u \neq w_{p}\) ; then \(u > w_{p}\) , whence \(u \geqslant w_{p+1}\) ; if \(u\) does not belong to \(X\) , then \(\alpha(w) < w_{p}\) , whence \(\alpha(w) < w_{p+1}\) ; then again \(u \in \mathrm{P}_{p+1}\) . Conversely, let \(u \in \mathrm{P}_{p+1}\) . We distinguish two cases:

\((\alpha)\) There exists no element \(v\) of \(\mathbf{M}(\mathbf{X})\) such that \(u = w_{p}v\) . By definition of \(\mathrm{P}_{p + 1}\) , \(u > w_{p}\) . Further, if \(u\notin \mathbf{X}\) , then \(\alpha (u)\neq w_p\) by the given hypothesis and \(\alpha (u) < w_{p + 1}\) since \(u\in \mathrm{P}_{p + 1}\) ; hence \(\alpha (u) < w_{p}\) . Hence \(u\in \mathrm{P}_p\) and \(u\neq w_{p}\) .

(β) There exists \(v\) in \(\mathbf{M}(\mathbf{X})\) such that \(u = w_{p}v\) . By Definition 2, of necessity, either \(w_{p} \in \mathbf{X}\) , \(v \in \mathbf{X}\) and \(w_{p} < v\) , or \(v \notin \mathbf{X}\) and \(\alpha(v) \leqslant w_{p} < v\) . In either case, \(v \in \mathbf{P}_{p+1}\) .

Then there exist an integer \(i \geqslant 0\) and an element \(w\) of \(\mathrm{M}(\mathrm{X})\) such that \(u = w_p^i w\) , and either \(w \in \mathrm{X}\) or \(w \notin \mathrm{X}\) and \(\alpha(w) \neq w_p\) . If \(i = 0\) , we have case \((\alpha)\) above, whence \(w \in \mathrm{P}_p\) and \(w \neq w_p\) . If \(i > 0\) , the proof of \((\beta)\) above establishes, by induction on \(i\) , the relations \(w \in \mathrm{P}_{p+1}\) and \(w \neq w_p\) . Suppose \(w \notin \mathrm{X}\) ; from \(w \in \mathrm{P}_{p+1}\) it follows that \(\alpha(w) \leqslant w_p\) and as \(\alpha(w) \neq w_p\) , we conclude that \(w \in \mathrm{P}_p\) . This completes the proof of (d).

Example. Suppose that X has two elements x, y; let X be ordered such that x < y. The construction given in the proof of Proposition 11 gives a set H with 14 elements of length \(\leqslant 5\) given in the following table:

\[
\begin{array}{l l l l} \mathrm{H} _ {1} & w _ {1} = x & w _ {2} = y \\ \mathrm{H} _ {2} & w _ {3} = (x y) \\ \mathrm{H} _ {3} & w _ {4} = (x (x y)) & w _ {5} = (y (x y)) \\ \mathrm{H} _ {4} & w _ {6} = (x (x (x y))) & w _ {7} = (y (x (x y))) & w _ {8} = (y (y (x y))) \\ \mathrm{H} _ {5} & w _ {9} = (x (x (x (x y)))) & w _ {10} = (y (x (x (x y)))) & w _ {11} = (y (y (x (x y)))) \\ & w _ {12} = (y (y (x y))) & w _ {13} = ((x y) (x (x y))) & w _ {14} = ((x y) (y (x y))), \end{array}
\]

(The elements of H have been numbered according to the total ordering chosen on each \(H_{n}\)) .

